\setcounter{equation}{0}
\section{Introduction}


Since the first computers, the trend has been to miniaturize these machines in order to increase their efficiency and convenience. But, clearly nature has limits to this miniaturization process. It was Benioff \cite{cps80} who proposed that the ultimate limit of miniaturization would be a quantum mechanical system. However, soon after Deutsch \cite{qtct85} showed that the only benefits of a quantum computer was not it's size and energy savings, a totally new way of computation was possible which made use of the entangled quantum particles and quantum parallelization. Following this, Shor \cite{ptafpf94} discovered a polynomial time prime factorization algorithm for a quantum computer which garnered a great interest into quantum computation. Together with Grover's algorithm for the fast searching of an unsorted database \cite{fqma96}, Shor's algorithm provided the motivation for the scientists to research the possibilities of realization of a quantum computing device.

Despite the rapid progress on the theory of quantum computation, the experiments for physical implementation of a quantum computer faced major challenges. Although very basic quantum computing devices have been accomplished using different techniques, realization of a useful quantum computer in the near future seems impossible.

Because of this reason quantum computer simulation by classical computers became a popular topic. Many libraries \footnote{Lib Quantum (\emph{http://www.enyo.de/libquantum/})} and standalone applications have been developed for this purpose, some of them free, \footnote {Jaquzzi (\emph{http://www.eng.buffalo.edu/~phygons/jaQuzzi/jaQuzzi.html}) } \footnote {Qudit (\emph{http://www.compsoc.nuigalway.ie/~damo642/QuantumSimulator/QuantumSimulator/QuantumQuditSimulator.htm})} some of them proprietary \footnote{Senko's simulator (\emph{http://www.senko-corp.co.jp/qcs/qcp/qcpe.html})},and some languages have been developed \footnote { QCL (\emph{http://tph.tuwien.ac.at/~oemer/qcl.html})} \footnote {Q language (\emph{http://sra.itc.it/people/serafini/qlang/})}. Because computation complexity increases exponentially with the number of qubits simulated, these attempts suffer great performance problems. To overcome these problems, this thesis discusses a method for quantum computer simulation on a distributed system. A similar approach has been tried by B. Oemer and Glendinning \cite{pgsq03}.
